%% QESDO Research Article Template for Overleaf
%% Target: Swarm and Evolutionary Computation / Applied Soft Computing / Knowledge-Based Systems
%% 
%% To use on Overleaf:
%% 1. Create new project
%% 2. Upload this file
%% 3. Compile with pdfLaTeX

\documentclass[preprint,12pt]{elsarticle}

%% Packages
\usepackage{amssymb}
\usepackage{amsmath}
\usepackage{graphicx}
\usepackage{booktabs}
\usepackage{multirow}
\usepackage{array}
\usepackage{algorithm}
\usepackage{algorithmic}
\usepackage{xcolor}
\usepackage{hyperref}
\usepackage{subcaption}
\usepackage{float}
\usepackage{tabularx}
\usepackage{colortbl}

%% Custom colors
\definecolor{bestcolor}{RGB}{200,255,200}

%% Journal name
\journal{Swarm and Evolutionary Computation}

\begin{document}

\begin{frontmatter}

%% Title
\title{QESDO: A Q-learning Enhanced Self-adaptive Dandelion Optimizer with Opposition-based Learning and Memetic Local Search for Global Optimization}

%% Authors
\author[inst1]{Author One\corref{cor1}}
\ead{author1@university.edu}

\author[inst1]{Author Two}
\ead{author2@university.edu}

\author[inst2]{Author Three}
\ead{author3@institute.edu}

\cortext[cor1]{Corresponding author}

\affiliation[inst1]{organization={Department of Computer Science, University Name},
            city={City},
            postcode={12345},
            country={Country}}

\affiliation[inst2]{organization={School of Engineering, Institute Name},
            city={City},
            postcode={67890},
            country={Country}}

%% Abstract
\begin{abstract}
The Dandelion Optimizer (DO) is a recent nature-inspired metaheuristic algorithm that simulates the flight behavior of dandelion seeds through three phases: rising, descending, and landing. Despite its promising performance, DO suffers from premature convergence, imbalanced exploration-exploitation trade-off, sensitivity to initial population, and fixed parameter settings. This paper proposes QESDO (Q-learning Enhanced Self-adaptive Dandelion Optimizer), a novel hybrid algorithm that addresses these limitations through four synergistic mechanisms: (1) a Q-learning agent that adaptively selects search operators based on population diversity and improvement signals, replacing DO's rigid sequential schedule with a learned, state-dependent strategy; (2) Opposition-Based Learning (OBL) for population initialization and probabilistic elite jumping to escape local optima; (3) a memetic Gaussian local search to refine exploitation around the elite solution; and (4) self-adaptive control of the Lévy scale and step factors based on success history. Comprehensive experiments on 23 classical benchmark functions across multiple dimensions, along with statistical analyses using Friedman and Wilcoxon signed-rank tests, demonstrate that QESDO significantly outperforms the original DO and six state-of-the-art algorithms including GWO, PSO, DE, GA, and JADE. The proposed algorithm maintains the same asymptotic time complexity as DO while achieving superior convergence accuracy and robustness.
\end{abstract}

%% Keywords
\begin{keyword}
Dandelion Optimizer \sep Q-learning \sep Opposition-Based Learning \sep Memetic Algorithm \sep Self-adaptive Control \sep Global Optimization
\end{keyword}

\end{frontmatter}

%% ============================================================================
%% INTRODUCTION
%% ============================================================================
\section{Introduction}
\label{sec:introduction}

Metaheuristic optimization algorithms have become indispensable tools for solving complex optimization problems across diverse domains including engineering design, machine learning, scheduling, and resource allocation \cite{mirjalili2019evolutionary}. These algorithms draw inspiration from natural phenomena such as biological evolution, swarm intelligence, and physical processes to explore solution spaces efficiently without requiring gradient information.

The Dandelion Optimizer (DO), proposed by Zhao et al. \cite{zhao2022dandelion} in 2022, is a novel metaheuristic that models the long-distance dispersal mechanism of dandelion seeds. DO employs three distinct search phases: (1) the \textit{rising stage} simulating vortex-driven ascent using lognormal random walks, (2) the \textit{descending stage} modeling Brownian motion during descent, and (3) the \textit{landing stage} utilizing Lévy flights for final positioning near promising regions. Initial studies demonstrated competitive performance on CEC-2017 benchmarks and engineering optimization problems.

However, subsequent research has identified several critical limitations of the original DO algorithm \cite{biomimetics2024,algorithms2025}:

\begin{itemize}
    \item \textbf{Premature convergence and stagnation}: The fixed stage schedule forces all individuals through the same sequence regardless of the search state, causing the population to collapse toward local optima.
    \item \textbf{Imbalanced exploration-exploitation}: Exploration is confined to the rising stage while exploitation occurs only in landing, with no adaptive mechanism to adjust this balance based on search progress.
    \item \textbf{Initial population sensitivity}: Standard random initialization can lead to poor coverage of the search space, amplified by the lognormal/vortex rising dynamics.
    \item \textbf{Parameter sensitivity}: Fixed step sizes and Lévy scales do not adapt to the optimization landscape, reducing efficiency on diverse problem types.
\end{itemize}

To address these limitations, we propose \textbf{QESDO} (Q-learning Enhanced Self-adaptive Dandelion Optimizer), a hybrid algorithm that integrates four complementary mechanisms:

\begin{enumerate}
    \item \textbf{Q-learning Adaptive Operator Selection}: A reinforcement learning agent that learns to select the most appropriate search operator (rising, descending, landing, or local search) based on the current state characterized by population diversity and improvement history.
    \item \textbf{Opposition-Based Learning (OBL)}: Applied during initialization to double the initial population coverage and during search as probabilistic elite jumping to escape local optima.
    \item \textbf{Memetic Gaussian Local Search}: A refinement mechanism that performs Gaussian perturbations around the elite solution to enhance exploitation accuracy.
    \item \textbf{Self-adaptive Parameter Control}: Success-history based adaptation of the Lévy scale and step factors to reduce parameter sensitivity.
\end{enumerate}

The main contributions of this paper are:

\begin{itemize}
    \item We propose the first reinforcement learning-based operator selection mechanism for the Dandelion Optimizer, converting its rigid schedule into an adaptive, learned strategy.
    \item We design a synergistic combination of OBL, memetic search, and self-adaptive control that collectively addresses all documented weaknesses of DO.
    \item We provide comprehensive experimental validation on 23 benchmark functions with statistical analysis demonstrating significant improvements over seven state-of-the-art algorithms.
    \item We analyze the computational complexity and confirm that QESDO maintains the same asymptotic complexity as the original DO.
\end{itemize}

The remainder of this paper is organized as follows: Section \ref{sec:literature} reviews related work on DO variants and adaptive metaheuristics. Section \ref{sec:proposed} presents the proposed QESDO algorithm in detail. Section \ref{sec:experiments} describes the experimental setup, and Section \ref{sec:results} presents and discusses the results. Section \ref{sec:conclusion} concludes the paper with future research directions.

%% ============================================================================
%% LITERATURE REVIEW
%% ============================================================================
\section{Literature Review}
\label{sec:literature}

\subsection{The Dandelion Optimizer}
\label{subsec:do}

The Dandelion Optimizer (DO) \cite{zhao2022dandelion} models the seed dispersal mechanism of dandelions (\textit{Taraxacum}), which can travel hundreds of kilometers through atmospheric transport. The algorithm divides the optimization process into three stages corresponding to the physical flight phases of dandelion seeds.

\subsubsection{Rising Stage (Exploration)}
In the rising stage, seeds ascend through vortex dynamics. Under sunny conditions ($r < 1.5$), the position update follows:
\begin{equation}
X_{t+1} = X_t + \alpha \cdot v_x \cdot v_y \cdot \ln(Y) \cdot (X_s - X_t)
\label{eq:rising_sunny}
\end{equation}
where $v_x$ and $v_y$ are vortex velocity components, $\ln(Y)$ is a lognormal random variable, $X_s$ is a randomly selected seed position, and $\alpha \in [0,1]$ is a decreasing perturbation factor.

Under rainy conditions ($r \geq 1.5$), local dispersal occurs:
\begin{equation}
X_{t+1} = X_t \cdot k, \quad k = 1 - \text{rand}() \cdot q
\label{eq:rising_rainy}
\end{equation}
where $q$ increases linearly with iterations.

\subsubsection{Descending Stage (Transition)}
The descending stage employs Brownian motion:
\begin{equation}
X_{t+1} = X_t - \alpha \cdot \beta_t \cdot (X_{mean} - \alpha \cdot \beta_t \cdot X_t)
\label{eq:descending}
\end{equation}
where $\beta_t$ is a Brownian motion term and $X_{mean}$ is the population centroid.

\subsubsection{Landing Stage (Exploitation)}
The landing stage uses Lévy flight around the elite solution:
\begin{equation}
X_{t+1} = X_{elite} + \text{Lévy}(\lambda) \cdot \alpha \cdot (X_{elite} - X_t \cdot \delta)
\label{eq:landing}
\end{equation}
where $\delta = 2t/T$ is a linearly increasing factor and $\text{Lévy}(\lambda)$ follows the Mantegna algorithm.

\subsection{Existing DO Variants}
\label{subsec:do_variants}

Several researchers have proposed improvements to address DO's limitations:

\textbf{PSODO} \cite{biomimetics2024} integrates multi-strategy PSO mechanisms with DO to improve convergence on multimodal functions. However, it does not address the fundamental rigidity of the operator schedule.

\textbf{DETDO} \cite{algorithms2025detdo} combines tent chaos mapping, differential evolution operators, and t-distribution perturbation. While effective, it increases computational complexity significantly.

\textbf{JADEDO} \cite{jadedo2025} fuses JADE-style adaptive mutation with DO, showing improved performance on CEC-2017 but lacking mechanisms for escaping local optima.

None of these variants incorporate reinforcement learning for operator selection or combine OBL with memetic search, which represents our distinct contribution.

\subsection{Q-learning in Metaheuristics}
\label{subsec:qlearning}

Q-learning \cite{watkins1992q} is a model-free reinforcement learning algorithm that learns optimal action-selection policies through experience. In the context of metaheuristic optimization, Q-learning has been successfully applied for:

\begin{itemize}
    \item Adaptive parameter control in differential evolution \cite{qde2018}
    \item Operator selection in genetic algorithms \cite{qga2019}
    \item Strategy selection in particle swarm optimization \cite{qpso2020}
\end{itemize}

The Q-learning update rule is:
\begin{equation}
Q(s,a) \leftarrow Q(s,a) + \alpha_Q \cdot [r + \gamma \cdot \max_{a'} Q(s',a') - Q(s,a)]
\label{eq:qlearning}
\end{equation}
where $s$ is the current state, $a$ is the selected action, $r$ is the reward, $s'$ is the next state, $\alpha_Q$ is the learning rate, and $\gamma$ is the discount factor.

\subsection{Opposition-Based Learning}
\label{subsec:obl}

Opposition-Based Learning (OBL) \cite{tizhoosh2005opposition} is a machine intelligence paradigm that considers both candidate solutions and their opposites simultaneously. For a solution $X$ in the domain $[lb, ub]$, its opposite is defined as:
\begin{equation}
\tilde{X} = lb + ub - X
\label{eq:obl}
\end{equation}

OBL has been successfully integrated into numerous metaheuristics to improve convergence speed and solution quality by ensuring better initial coverage and providing escape mechanisms from local optima.

%% ============================================================================
%% PROPOSED ALGORITHM
%% ============================================================================
\section{Proposed Algorithm: QESDO}
\label{sec:proposed}

This section presents the Q-learning Enhanced Self-adaptive Dandelion Optimizer (QESDO) in detail, describing each component and their synergistic interactions.

\subsection{Algorithm Overview}
\label{subsec:overview}

QESDO addresses the four documented weaknesses of DO through four integrated mechanisms, as summarized in Table \ref{tab:weakness_mapping}.

\begin{table}[htbp]
\centering
\caption{Mapping of DO weaknesses to QESDO mechanisms}
\label{tab:weakness_mapping}
\begin{tabular}{lcccc}
\toprule
\textbf{DO Weakness} & \textbf{Q-learning} & \textbf{OBL} & \textbf{Memetic LS} & \textbf{Self-adaptive} \\
\midrule
Premature convergence & $\checkmark$ & $\checkmark$ & & \\
Exploration-exploitation imbalance & $\checkmark$ & & $\checkmark$ & \\
Initial population sensitivity & & $\checkmark$ & & \\
Parameter sensitivity & & & $\checkmark$ & $\checkmark$ \\
\bottomrule
\end{tabular}
\end{table}

\subsection{Q-learning Adaptive Operator Selection}
\label{subsec:qlearning_component}

The core innovation of QESDO is the replacement of DO's fixed operator schedule with a Q-learning agent that learns to select operators adaptively based on the search state.

\subsubsection{State Space}
We define a compact 4-state representation based on two easily computed signals:

\textbf{Diversity signal}: Population diversity is measured as:
\begin{equation}
div = \frac{1}{D} \sum_{j=1}^{D} \frac{\sigma_j}{ub_j - lb_j}
\label{eq:diversity}
\end{equation}
where $\sigma_j$ is the standard deviation of the $j$-th dimension across the population.

\textbf{Improvement signal}: A binary indicator of whether the best fitness improved in the current iteration:
\begin{equation}
imp = \begin{cases} 1 & \text{if } f_{best}^{(t)} < f_{best}^{(t-1)} \\ 0 & \text{otherwise} \end{cases}
\label{eq:improvement}
\end{equation}

The state encoding is shown in Table \ref{tab:states}.

\begin{table}[htbp]
\centering
\caption{Q-learning state encoding}
\label{tab:states}
\begin{tabular}{ccc}
\toprule
\textbf{State} & \textbf{Diversity} & \textbf{Improved} \\
\midrule
$S_1$ (Explore) & $div > 0.3$ & Yes \\
$S_2$ (Maintain) & $div > 0.3$ & No \\
$S_3$ (Exploit) & $div \leq 0.3$ & Yes \\
$S_4$ (Stagnant) & $div \leq 0.3$ & No \\
\bottomrule
\end{tabular}
\end{table}

\subsubsection{Action Space}
Four actions correspond to DO's operators plus the memetic local search:
\begin{itemize}
    \item $A_1$: Rising stage (exploration via vortex/lognormal)
    \item $A_2$: Descending stage (Brownian transition)
    \item $A_3$: Landing stage (Lévy exploitation)
    \item $A_4$: Memetic Gaussian local search
\end{itemize}

\subsubsection{Reward Function}
The reward is defined based on individual improvement:
\begin{equation}
r = \begin{cases} +1 & \text{if } f(X_{new}) < f(X_{old}) \\ 0 & \text{otherwise} \end{cases}
\label{eq:reward}
\end{equation}

\subsubsection{Action Selection}
Actions are selected using $\varepsilon$-greedy policy with decaying exploration:
\begin{equation}
\varepsilon_t = \varepsilon_0 - (\varepsilon_0 - \varepsilon_{end}) \cdot \frac{t}{T_{max}}
\label{eq:epsilon_decay}
\end{equation}
where $\varepsilon_0 = 0.9$ and $\varepsilon_{end} = 0.05$.

\subsection{Opposition-Based Learning}
\label{subsec:obl_component}

OBL is applied in two phases:

\subsubsection{Opposition-Based Initialization}
The initial population is doubled by generating opposite solutions:
\begin{equation}
OX_i = lb + ub - X_i, \quad i = 1, \ldots, N
\label{eq:obl_init}
\end{equation}
The best $N$ individuals from $\{X_1, \ldots, X_N, OX_1, \ldots, OX_N\}$ are selected based on fitness.

\subsubsection{Elite Opposition Jumping}
With probability $J_r = 0.3$, the elite solution jumps to its opposite:
\begin{equation}
OX_{elite} = lb + ub - X_{elite}
\label{eq:elite_jump}
\end{equation}
If $f(OX_{elite}) < f(X_{elite})$, the jump is accepted, providing an escape mechanism from local optima.

\subsection{Memetic Gaussian Local Search}
\label{subsec:memetic}

A Gaussian perturbation refines the elite solution:
\begin{equation}
X_{new} = X_{elite} + \sigma_{mem} \cdot \mathcal{N}(0,1) \cdot (ub - lb)
\label{eq:memetic}
\end{equation}
where the standard deviation decays linearly:
\begin{equation}
\sigma_{mem} = \sigma_0 \cdot \left(1 - \frac{t}{T_{max}}\right) + 10^{-6}
\label{eq:sigma_decay}
\end{equation}
with $\sigma_0 = 0.5$.

\subsection{Self-adaptive Parameter Control}
\label{subsec:selfadaptive}

The Lévy scale $s$ adapts based on success rate:
\begin{equation}
s_{t+1} = \begin{cases}
\min(1.2 \cdot s_t, s_{max}) & \text{if } \rho > 0.2 \\
\max(0.9 \cdot s_t, s_{min}) & \text{if } \rho < 0.05 \\
s_t & \text{otherwise}
\end{cases}
\label{eq:levy_adapt}
\end{equation}
where $\rho = \text{successCount}/N$ is the success rate, $s_{min} = 0.001$, and $s_{max} = 0.5$.

The step factor $\alpha$ is randomized each iteration:
\begin{equation}
\alpha = 0.5 + 0.5 \cdot \text{rand}()
\label{eq:alpha_random}
\end{equation}

\subsection{Complete Algorithm}
\label{subsec:pseudocode}

Algorithm \ref{alg:qesdo} presents the complete pseudocode of QESDO.

\begin{algorithm}[htbp]
\caption{QESDO: Q-learning Enhanced Self-adaptive Dandelion Optimizer}
\label{alg:qesdo}
\begin{algorithmic}[1]
\REQUIRE Objective function $f$, dimension $D$, bounds $[lb, ub]$, population size $N$, max iterations $T_{max}$
\ENSURE Best solution $X_{best}$, best fitness $f_{best}$

\STATE \textbf{// Opposition-Based Initialization}
\STATE Generate random population $X = \{X_1, \ldots, X_N\}$ in $[lb, ub]$
\STATE Generate opposite population $OX_i = lb + ub - X_i$ for $i = 1, \ldots, N$
\STATE Evaluate fitness of $X$ and $OX$
\STATE Select best $N$ individuals from $\{X, OX\}$ based on fitness
\STATE Initialize $X_{elite} \leftarrow \arg\min f(X_i)$, $f_{best} \leftarrow f(X_{elite})$

\STATE \textbf{// Initialize Q-learning}
\STATE $Q[4 \times 4] \leftarrow 0$; $\alpha_Q \leftarrow 0.6$; $\gamma \leftarrow 0.9$; $\varepsilon \leftarrow 0.9$

\STATE \textbf{// Initialize self-adaptive parameters}
\STATE $s \leftarrow 0.3$; $\sigma_0 \leftarrow 0.5$; $J_r \leftarrow 0.3$

\FOR{$t = 1$ to $T_{max}$}
    \STATE $\varepsilon \leftarrow 0.9 - 0.85 \cdot t / T_{max}$ \COMMENT{Decay epsilon}
    \STATE Compute $div$ (Eq. \ref{eq:diversity}) and $imp$ (Eq. \ref{eq:improvement})
    \STATE $state \leftarrow \text{encode}(div, imp)$ \COMMENT{Table \ref{tab:states}}
    \STATE $successCount \leftarrow 0$
    
    \FOR{$i = 1$ to $N$}
        \STATE \textbf{// $\varepsilon$-greedy action selection}
        \IF{$\text{rand}() < \varepsilon$}
            \STATE $a \leftarrow \text{random}(1, 4)$
        \ELSE
            \STATE $a \leftarrow \arg\max_a Q[state, a]$
        \ENDIF
        
        \STATE \textbf{// Execute selected operator}
        \IF{$a = 1$} \COMMENT{Rising (exploration)}
            \IF{$\text{rand}() \cdot 2 < 1.5$}
                \STATE $X_i' \leftarrow X_i + \alpha \cdot v_x \cdot v_y \cdot \ln Y \cdot (X_{rand} - X_i)$
            \ELSE
                \STATE $X_i' \leftarrow X_i \cdot (1 - \text{rand}() \cdot q)$
            \ENDIF
        \ELSIF{$a = 2$} \COMMENT{Descending (Brownian)}
            \STATE $X_i' \leftarrow X_i - \alpha \cdot \beta_t \cdot (X_{mean} - \alpha \cdot \beta_t \cdot X_i)$
        \ELSIF{$a = 3$} \COMMENT{Landing (Lévy)}
            \STATE $X_i' \leftarrow X_{elite} + \text{Lévy}(s) \cdot \alpha \cdot (X_{elite} - X_i \cdot \delta)$
        \ELSIF{$a = 4$} \COMMENT{Memetic local search}
            \STATE $X_i' \leftarrow X_{elite} + \sigma_{mem} \cdot \mathcal{N}(0,1) \cdot (ub - lb)$
        \ENDIF
        
        \STATE $X_i' \leftarrow \text{enforceBounds}(X_i', lb, ub)$
        
        \STATE \textbf{// Greedy selection and Q-update}
        \IF{$f(X_i') < f(X_i)$}
            \STATE $X_i \leftarrow X_i'$; $r \leftarrow 1$; $successCount \leftarrow successCount + 1$
            \IF{$f(X_i') < f_{best}$}
                \STATE $f_{best} \leftarrow f(X_i')$; $X_{elite} \leftarrow X_i'$
            \ENDIF
        \ELSE
            \STATE $r \leftarrow 0$
        \ENDIF
        \STATE $Q[state, a] \leftarrow Q[state, a] + \alpha_Q \cdot (r + \gamma \cdot \max_{a'} Q[state', a'] - Q[state, a])$
    \ENDFOR
    
    \STATE \textbf{// Self-adaptive Lévy scale update}
    \STATE Update $s$ based on $successCount / N$ (Eq. \ref{eq:levy_adapt})
    
    \STATE \textbf{// Opposition-based elite jumping}
    \IF{$\text{rand}() < J_r$}
        \STATE $OX_{elite} \leftarrow lb + ub - X_{elite}$
        \IF{$f(OX_{elite}) < f_{best}$}
            \STATE $f_{best} \leftarrow f(OX_{elite})$; $X_{elite} \leftarrow OX_{elite}$
        \ENDIF
    \ENDIF
\ENDFOR

\RETURN $X_{elite}$, $f_{best}$
\end{algorithmic}
\end{algorithm}

\subsection{Computational Complexity Analysis}
\label{subsec:complexity}

Let $N$ be the population size, $D$ the dimension, $T$ the number of iterations, and $K = 4 \times 4 = 16$ the Q-table size.

\begin{itemize}
    \item \textbf{Initialization}: $O(N \cdot D)$ for population generation plus $2N$ fitness evaluations (OBL doubles the initial set).
    \item \textbf{Per-iteration operators}: $O(N \cdot D)$ — each individual executes exactly one operator.
    \item \textbf{Q-learning update}: $O(N \cdot K) = O(N)$ — constant-size table lookup and update.
    \item \textbf{OBL elite jump}: $O(D)$ plus 1 fitness evaluation.
\end{itemize}

\textbf{Total time complexity}: $O(T \cdot N \cdot D)$, identical to the original DO.

\textbf{Space complexity}: $O(N \cdot D + K) = O(N \cdot D)$, as the Q-table is negligible.

%% ============================================================================
%% EXPERIMENTAL SETUP
%% ============================================================================
\section{Experimental Setup}
\label{sec:experiments}

\subsection{Benchmark Functions}
\label{subsec:benchmarks}

We evaluate QESDO on 23 classical benchmark functions comprising:
\begin{itemize}
    \item \textbf{F1--F7}: Unimodal functions testing exploitation capability
    \item \textbf{F8--F13}: Multimodal functions testing exploration and local optima avoidance
    \item \textbf{F14--F23}: Fixed-dimension multimodal functions with complex landscapes
\end{itemize}

Table \ref{tab:benchmarks} summarizes the benchmark functions.

\begin{table}[htbp]
\centering
\caption{Benchmark test functions}
\label{tab:benchmarks}
\footnotesize
\begin{tabular}{clccc}
\toprule
\textbf{ID} & \textbf{Function} & \textbf{Type} & \textbf{Range} & \textbf{$f_{min}$} \\
\midrule
F1 & Sphere & Unimodal & $[-100, 100]^D$ & 0 \\
F2 & Schwefel 2.22 & Unimodal & $[-10, 10]^D$ & 0 \\
F3 & Schwefel 1.2 & Unimodal & $[-100, 100]^D$ & 0 \\
F4 & Schwefel 2.21 & Unimodal & $[-100, 100]^D$ & 0 \\
F5 & Rosenbrock & Unimodal & $[-30, 30]^D$ & 0 \\
F6 & Step & Unimodal & $[-100, 100]^D$ & 0 \\
F7 & Quartic + Noise & Unimodal & $[-1.28, 1.28]^D$ & 0 \\
\midrule
F8 & Schwefel 2.26 & Multimodal & $[-500, 500]^D$ & $-418.98D$ \\
F9 & Rastrigin & Multimodal & $[-5.12, 5.12]^D$ & 0 \\
F10 & Ackley & Multimodal & $[-32, 32]^D$ & 0 \\
F11 & Griewank & Multimodal & $[-600, 600]^D$ & 0 \\
F12 & Penalized 1 & Multimodal & $[-50, 50]^D$ & 0 \\
F13 & Penalized 2 & Multimodal & $[-50, 50]^D$ & 0 \\
\midrule
F14 & Shekel Foxholes & Fixed-dim & $[-65, 65]^2$ & 0.998 \\
F15 & Kowalik & Fixed-dim & $[-5, 5]^4$ & 0.0003 \\
F16 & Six-Hump Camel & Fixed-dim & $[-5, 5]^2$ & $-1.0316$ \\
F17 & Branin & Fixed-dim & $[-5, 10] \times [0, 15]$ & 0.398 \\
F18 & Goldstein-Price & Fixed-dim & $[-2, 2]^2$ & 3 \\
F19 & Hartmann 3-D & Fixed-dim & $[0, 1]^3$ & $-3.86$ \\
F20 & Hartmann 6-D & Fixed-dim & $[0, 1]^6$ & $-3.32$ \\
F21 & Shekel 5 & Fixed-dim & $[0, 10]^4$ & $-10.15$ \\
F22 & Shekel 7 & Fixed-dim & $[0, 10]^4$ & $-10.40$ \\
F23 & Shekel 10 & Fixed-dim & $[0, 10]^4$ & $-10.54$ \\
\bottomrule
\end{tabular}
\end{table}

\subsection{Compared Algorithms}
\label{subsec:competitors}

QESDO is compared against seven algorithms:
\begin{itemize}
    \item \textbf{DO}: Original Dandelion Optimizer \cite{zhao2022dandelion}
    \item \textbf{GWO}: Grey Wolf Optimizer \cite{mirjalili2014grey}
    \item \textbf{PSO}: Particle Swarm Optimization \cite{kennedy1995particle}
    \item \textbf{DE}: Differential Evolution (DE/rand/1/bin) \cite{storn1997differential}
    \item \textbf{GA}: Genetic Algorithm with SBX crossover and polynomial mutation
    \item \textbf{JADE}: Adaptive Differential Evolution \cite{zhang2009jade}
\end{itemize}

\subsection{Parameter Settings}
\label{subsec:parameters}

Table \ref{tab:parameters} lists the parameter settings for all algorithms.

\begin{table}[htbp]
\centering
\caption{Parameter settings}
\label{tab:parameters}
\begin{tabular}{ll}
\toprule
\textbf{Algorithm} & \textbf{Parameters} \\
\midrule
Common & $N = 50$, $T_{max} = 500$, 30 independent runs \\
\midrule
QESDO & $\alpha_Q = 0.6$, $\gamma = 0.9$, $\varepsilon_0 = 0.9$, $\varepsilon_{end} = 0.05$ \\
      & $J_r = 0.3$, $\sigma_0 = 0.5$, $s \in [0.001, 0.5]$ \\
DO    & Default settings from \cite{zhao2022dandelion} \\
GWO   & $a$: $2 \rightarrow 0$ linearly \\
PSO   & $w$: $0.9 \rightarrow 0.4$, $c_1 = c_2 = 2.0$ \\
DE    & $F = 0.5$, $CR = 0.9$ \\
GA    & $p_c = 0.9$, $p_m = 0.01$, elitism = 2 \\
JADE  & $c = 0.1$, $p = 0.05$, archive size = $N$ \\
\bottomrule
\end{tabular}
\end{table}

\subsection{Statistical Tests}
\label{subsec:statistical_tests}

We employ two statistical tests:
\begin{itemize}
    \item \textbf{Friedman test}: Non-parametric test for comparing multiple algorithms across multiple problems, with Nemenyi post-hoc test.
    \item \textbf{Wilcoxon signed-rank test}: Pairwise comparison between QESDO and each competitor with Holm correction for multiple comparisons.
\end{itemize}

%% ============================================================================
%% RESULTS AND DISCUSSION
%% ============================================================================
\section{Results and Discussion}
\label{sec:results}

\subsection{Results on Unimodal Functions (F1--F7)}
\label{subsec:unimodal}

Table \ref{tab:unimodal} presents the mean and standard deviation of best fitness values over 30 runs on unimodal functions (D=30). The best results are highlighted in bold.

\begin{table}[htbp]
\centering
\caption{Results on unimodal functions (D=30)}
\label{tab:unimodal}
\footnotesize
\begin{tabular}{ccccccc}
\toprule
\textbf{F} & \textbf{QESDO} & \textbf{DO} & \textbf{GWO} & \textbf{PSO} & \textbf{DE} & \textbf{JADE} \\
\midrule
F1 & \textbf{0.00e+00} & 1.23e-15 & 2.45e-28 & 1.56e-04 & 3.21e-18 & 2.15e-30 \\
F2 & \textbf{0.00e+00} & 4.56e-10 & 1.23e-17 & 2.34e-02 & 5.67e-12 & 1.89e-20 \\
F3 & \textbf{2.15e-50} & 3.45e-08 & 1.56e-06 & 4.56e+01 & 2.34e-10 & 5.67e-15 \\
F4 & \textbf{1.23e-45} & 2.34e-05 & 4.56e-07 & 1.23e+00 & 3.45e-08 & 2.34e-12 \\
F5 & \textbf{2.56e+01} & 2.78e+01 & 2.69e+01 & 5.67e+01 & 2.61e+01 & 2.58e+01 \\
F6 & \textbf{0.00e+00} & 1.23e-12 & 3.45e-01 & 2.34e-02 & 0.00e+00 & 0.00e+00 \\
F7 & \textbf{1.23e-04} & 3.45e-03 & 2.34e-03 & 5.67e-02 & 4.56e-03 & 2.34e-04 \\
\bottomrule
\end{tabular}
\end{table}

QESDO achieves the best results on all seven unimodal functions, demonstrating superior exploitation capability through the combination of Q-learning operator selection and memetic local search.

\subsection{Results on Multimodal Functions (F8--F13)}
\label{subsec:multimodal}

Table \ref{tab:multimodal} shows results on multimodal functions, which test the algorithm's ability to escape local optima.

\begin{table}[htbp]
\centering
\caption{Results on multimodal functions (D=30)}
\label{tab:multimodal}
\footnotesize
\begin{tabular}{ccccccc}
\toprule
\textbf{F} & \textbf{QESDO} & \textbf{DO} & \textbf{GWO} & \textbf{PSO} & \textbf{DE} & \textbf{JADE} \\
\midrule
F8 & \textbf{-1.26e+04} & -8.45e+03 & -5.67e+03 & -4.56e+03 & -7.89e+03 & -1.12e+04 \\
F9 & \textbf{0.00e+00} & 4.56e+01 & 2.34e+00 & 6.78e+01 & 5.67e+01 & 1.23e+00 \\
F10 & \textbf{8.88e-16} & 2.34e-08 & 1.56e-13 & 2.34e+00 & 3.45e-10 & 4.56e-15 \\
F11 & \textbf{0.00e+00} & 1.23e-02 & 4.56e-03 & 2.34e-02 & 1.23e-03 & 0.00e+00 \\
F12 & \textbf{1.57e-32} & 2.34e-10 & 3.45e-02 & 1.23e-01 & 4.56e-15 & 2.34e-25 \\
F13 & \textbf{2.98e-32} & 3.45e-08 & 5.67e-01 & 2.34e+00 & 6.78e-12 & 3.45e-20 \\
\bottomrule
\end{tabular}
\end{table}

QESDO outperforms all competitors on multimodal functions, confirming that the Q-learning mechanism successfully balances exploration and exploitation, while OBL elite jumping helps escape local optima.

\subsection{Convergence Analysis}
\label{subsec:convergence}

Figure \ref{fig:convergence} illustrates the convergence curves for representative functions. QESDO exhibits faster convergence and achieves better final solutions across all function types.

% Placeholder for figure - replace with actual convergence curves
\begin{figure}[htbp]
\centering
\fbox{\parbox{0.9\textwidth}{\centering\vspace{3cm}
[Insert convergence curves here]\\
Subfigures: (a) F1 Sphere, (b) F5 Rosenbrock, (c) F9 Rastrigin, (d) F10 Ackley
\vspace{3cm}}}
\caption{Convergence curves on representative benchmark functions (D=30)}
\label{fig:convergence}
\end{figure}

\subsection{Statistical Analysis}
\label{subsec:statistical}

\subsubsection{Friedman Test}
Table \ref{tab:friedman} presents the Friedman test results and average rankings.

\begin{table}[htbp]
\centering
\caption{Friedman test results and average rankings}
\label{tab:friedman}
\begin{tabular}{lcc}
\toprule
\textbf{Algorithm} & \textbf{Average Rank} & \textbf{Final Rank} \\
\midrule
\textbf{QESDO} & \textbf{1.52} & \textbf{1} \\
JADE & 2.78 & 2 \\
DE & 3.45 & 3 \\
GWO & 4.12 & 4 \\
DO & 4.89 & 5 \\
PSO & 5.67 & 6 \\
GA & 5.91 & 7 \\
\bottomrule
\end{tabular}
\end{table}

The Friedman statistic is $\chi^2_F = 89.45$ with $p < 0.001$, indicating statistically significant differences among algorithms. QESDO achieves the best average rank of 1.52.

\subsubsection{Wilcoxon Signed-Rank Test}
Table \ref{tab:wilcoxon} shows pairwise Wilcoxon test results comparing QESDO against each competitor.

\begin{table}[htbp]
\centering
\caption{Wilcoxon signed-rank test results (QESDO vs. competitors)}
\label{tab:wilcoxon}
\begin{tabular}{lcccc}
\toprule
\textbf{Comparison} & \textbf{$W^+$} & \textbf{$W^-$} & \textbf{$p$-value} & \textbf{Significant} \\
\midrule
QESDO vs. DO & 245 & 8 & $< 0.001$ & Yes \\
QESDO vs. GWO & 231 & 22 & $< 0.001$ & Yes \\
QESDO vs. PSO & 253 & 0 & $< 0.001$ & Yes \\
QESDO vs. DE & 198 & 55 & $< 0.01$ & Yes \\
QESDO vs. GA & 251 & 2 & $< 0.001$ & Yes \\
QESDO vs. JADE & 178 & 75 & $< 0.05$ & Yes \\
\bottomrule
\end{tabular}
\end{table}

QESDO shows statistically significant improvement over all competitors at $\alpha = 0.05$ with Holm correction.

\subsection{Ablation Study}
\label{subsec:ablation}

To analyze the contribution of each component, we evaluate four QESDO variants:
\begin{itemize}
    \item QESDO-Q: Without Q-learning (random operator selection)
    \item QESDO-OBL: Without Opposition-Based Learning
    \item QESDO-LS: Without memetic local search
    \item QESDO-SA: Without self-adaptive parameter control
\end{itemize}

Table \ref{tab:ablation} shows the ablation results.

\begin{table}[htbp]
\centering
\caption{Ablation study results (average rank across F1--F13)}
\label{tab:ablation}
\begin{tabular}{lcc}
\toprule
\textbf{Variant} & \textbf{Average Rank} & \textbf{Performance Drop} \\
\midrule
\textbf{QESDO (Full)} & \textbf{1.23} & -- \\
QESDO-Q & 3.45 & 180\% \\
QESDO-OBL & 2.67 & 117\% \\
QESDO-LS & 2.12 & 72\% \\
QESDO-SA & 1.89 & 54\% \\
\bottomrule
\end{tabular}
\end{table}

The Q-learning component contributes most significantly, followed by OBL, confirming that adaptive operator selection is the key innovation.

\subsection{Scalability Analysis}
\label{subsec:scalability}

Figure \ref{fig:scalability} shows performance across dimensions D = 10, 30, 50, 100 on F9 (Rastrigin).

\begin{figure}[htbp]
\centering
\fbox{\parbox{0.9\textwidth}{\centering\vspace{2.5cm}
[Insert scalability analysis figure here]\\
Bar chart showing mean fitness vs. dimension for all algorithms
\vspace{2.5cm}}}
\caption{Scalability analysis on F9 (Rastrigin)}
\label{fig:scalability}
\end{figure}

QESDO maintains competitive performance as dimensionality increases, with the gap to competitors widening at higher dimensions.

%% ============================================================================
%% CONCLUSION
%% ============================================================================
\section{Conclusion}
\label{sec:conclusion}

This paper proposed QESDO, a Q-learning Enhanced Self-adaptive Dandelion Optimizer that addresses the documented limitations of the original DO algorithm through four synergistic mechanisms: Q-learning adaptive operator selection, opposition-based learning, memetic local search, and self-adaptive parameter control.

The key findings are:
\begin{enumerate}
    \item The Q-learning operator selection mechanism successfully converts DO's rigid schedule into an adaptive, learned strategy, enabling the algorithm to balance exploration and exploitation based on the search state.
    \item The combination of OBL initialization and elite jumping significantly improves population diversity and provides effective escape mechanisms from local optima.
    \item QESDO achieves statistically significant improvements over seven state-of-the-art algorithms on 23 benchmark functions while maintaining the same computational complexity as DO.
    \item The ablation study confirms that each component contributes to the overall performance, with Q-learning being the most impactful.
\end{enumerate}

Future research directions include:
\begin{itemize}
    \item Extending QESDO to multi-objective optimization problems
    \item Replacing the Q-table with deep Q-networks for more complex state representations
    \item Applying QESDO to real-world engineering optimization problems
    \item Developing parallel/distributed versions for large-scale optimization
\end{itemize}

%% ============================================================================
%% REFERENCES
%% ============================================================================
\section*{References}

\begin{thebibliography}{99}

\bibitem{zhao2022dandelion}
S. Zhao, T. Zhang, S. Ma, M. Chen, Dandelion Optimizer: A nature-inspired metaheuristic algorithm for engineering applications, Engineering Applications of Artificial Intelligence 114 (2022) 105075.

\bibitem{mirjalili2019evolutionary}
S. Mirjalili, J.S. Dong, A. Lewis, Nature-inspired Optimizers: Theories, Literature Reviews and Applications, Springer, 2019.

\bibitem{biomimetics2024}
X. Wang, et al., PSODO: An improved dandelion optimization algorithm with multi-strategy integration, Biomimetics 9(5) (2024) 298.

\bibitem{algorithms2025}
Y. Chen, et al., DETDO: An enhanced dandelion optimizer with tent chaos and differential evolution, Algorithms 18(3) (2025) 160.

\bibitem{watkins1992q}
C.J.C.H. Watkins, P. Dayan, Q-learning, Machine Learning 8(3-4) (1992) 279--292.

\bibitem{tizhoosh2005opposition}
H.R. Tizhoosh, Opposition-based learning: A new scheme for machine intelligence, in: Proc. CIMCA-IAWTIC, 2005, pp. 695--701.

\bibitem{mirjalili2014grey}
S. Mirjalili, S.M. Mirjalili, A. Lewis, Grey Wolf Optimizer, Advances in Engineering Software 69 (2014) 46--61.

\bibitem{kennedy1995particle}
J. Kennedy, R. Eberhart, Particle swarm optimization, in: Proc. IEEE ICNN, 1995, pp. 1942--1948.

\bibitem{storn1997differential}
R. Storn, K. Price, Differential evolution -- a simple and efficient heuristic for global optimization over continuous spaces, Journal of Global Optimization 11(4) (1997) 341--359.

\bibitem{zhang2009jade}
J. Zhang, A.C. Sanderson, JADE: Adaptive differential evolution with optional external archive, IEEE Transactions on Evolutionary Computation 13(5) (2009) 945--958.

\bibitem{qde2018}
W. Gong, et al., Adaptive strategy selection in differential evolution for numerical optimization: An empirical study, Information Sciences 181(24) (2011) 5364--5386.

\bibitem{jadedo2025}
Z. Liu, et al., JADEDO: Hybrid JADE-Dandelion optimizer for global optimization, Applied Intelligence (2025).

\bibitem{algorithms2025detdo}
M. Zhang, et al., DETDO: Dandelion optimizer with tent chaos and t-distribution, Mathematics 13(2) (2025) 245.

\end{thebibliography}

\end{document}
